% Source-derived chapter 02: Terminology
% Original source: bounds-stalks-perverse-sheaves-characteristic-p-shende
\chapter{Terminology}
\label{bounds-stalks-perverse-sheaves-characteristic-p-shende-chapter-02}
\source{bounds-stalks-perverse-sheaves-characteristic-p-shende}

\begin{remark}[Source section]
\label{bounds-stalks-perverse-sheaves-characteristic-p-shende-chapter-02-source}
\source{bounds-stalks-perverse-sheaves-characteristic-p-shende}
\source{bounds-stalks-perverse-sheaves-characteristic-p-shende}
This chapter reproduces the corresponding section of the retrieved
LaTeX source, including its definitions, statements, examples, and
proofs. It has not yet been translated into Lean.
\notready
\end{remark}

% The source section title was: Terminology
We review some notation and terminology from \cite{Beilinson} and \cite{saito1}. (Our formulations of the definitions are mainly adapted from \cite{saito1}). All schemes are over a perfect field $k$, which in the application we can specialize to be the algebraic closure of a finite field.

\begin{definition}
\source{bounds-stalks-perverse-sheaves-characteristic-p-shende}
\notready\label{C-transversal-1} \cite[Definition 3.5(1)]{saito1} Let $X$ be a smooth scheme over $k$ and let $C \subseteq T^* X$ be a closed conical subset of the cotangent bundle. Let $f: X\to Y$ be a morphism of smooth schemes over $k$.
\source{bounds-stalks-perverse-sheaves-characteristic-p-shende}

We say that $f: X \to Y$ is \emph{$C$-transversal} if the inverse image $df^{-1}(C)$ by the canonical morphism $X \times_Y T^* Y \to T^* X$ is a subset of the zero-section $X \subseteq X \times_Y T^* Y$.
\end{definition}

\begin{definition}
\source{bounds-stalks-perverse-sheaves-characteristic-p-shende}
\notready\cite[(1.2)]{Beilinson}\label{circ-forward} In the same setting as Definition \ref{C-transversal-1}, if $f$ is proper, let $f_\circ C$ be pushforward from $X \times_Y T^* Y$ to $T^*Y $ of the the inverse image $df^{-1}(C)$.
\source{bounds-stalks-perverse-sheaves-characteristic-p-shende}

\end{definition}

\begin{definition}
\source{bounds-stalks-perverse-sheaves-characteristic-p-shende}
\notready\cite[(2.3)]{saito-direct}\label{shriek-forward}  In the same setting as Definitions \ref{C-transversal-1} and \ref{circ-forward}, let $A$ be an algebraic cycle of codimension $\dim X$ supported on $C$. Assume also that $f_\circ C$ has dimension $\dim Y$.
\source{bounds-stalks-perverse-sheaves-characteristic-p-shende}

Let $f_!  A$ be the pushforward from $X \times_Y T^* Y$ to $T^*Y $ of the intersection-theoretic inverse image $df^* C$. \end{definition} 



\begin{definition}
\source{bounds-stalks-perverse-sheaves-characteristic-p-shende}
\notready\cite[Definition 3.1]{saito1} Let $X$ be a smooth scheme over $k$ and let $C \subseteq T^* X$ be a closed conical subset of the cotangent bundle. Let $h: W \to X$ be a morphism of smooth schemes over $k$.

Let $h^* C$ be the pullback of $C$ from $T^* X$ to $W \times_X T^* X$ and let $K$ be the inverse image of the $0$-section $W \subseteq T^* W$ by the canonical morphism $dh: W \times_X T^* X \to T^* W$. 

We say that $h: W\to X$ is \emph{$C$-transversal} if the intersection $h^* C \cap K$ is a subset of the zero-section $W \subseteq W \times_X T^* X$.

If $h: W \to X$ is $C$-transversal, we define a closed conical subset $h^\circ C \subseteq T^* W$ as the image of $h^* C$ under $dh$ (it is closed by \cite[Lemma 3.1]{saito1}). \end{definition}

\begin{definition}
\source{bounds-stalks-perverse-sheaves-characteristic-p-shende}
\notready\cite[Definition 3.5(2)]{saito1} We say that a pair of morphisms $h: W \to X$ and $f:W\to Y$ of smooth schemes over $k$ is \emph{$C$-transversal}, for $C \subseteq T^* X$ a closed conical subset of the cotangent bundle, if $h$ is $C$-transversal and $f$ is $h^\circ C$-transversal. \end{definition}

\begin{definition}
\source{bounds-stalks-perverse-sheaves-characteristic-p-shende}
\notready\cite[Th. Finitude, Definition 2.12]{sga4h} We say that a morphism $f: W \to Y$ is locally acyclic relative to $K \in D^b_x(W, \mathbb F_\ell)$ if for each geometric point $x \in W$ geometric point  $t \in Y$ specializing to $f(x)$, $W_x$ the henselization of $W$ at $x$ and $W_{x,t}$ the fiber of $X_x$ over $t$, the natural map $H^* ( W_x, K) \to H^* ( W_{x,t}, K)$ is an isomorphism.  \end{definition}

For us, the main advantage of local acyclicity is that, when $Y$ is a smooth curve, the local acyclicity of $f$ implies that $R\Phi_f K$ vanishes, since taking $t$ the generic point, $H^* ( W_{x,t}, K)$ is the stalk of $R \Psi_ K$ at $x$ and $H^* ( W_x, K) $ is the stalk of $K$ at $x$, so the map is an isomorphism if and only if the mapping cone $R(\Phi_f K)_x$ vanishes.

\begin{definition}
\source{bounds-stalks-perverse-sheaves-characteristic-p-shende}
\notready\cite[1.3]{Beilinson} For $K \in D^b_c(X, \mathbb F_\ell)$, let the \emph{singular support $SS(K)$ of $K$} be the smallest closed conical subset $C \in T^* X$ such that for every $C$-transversal pair $h: W \to X$ and $f: W\to Y$, the morphism $f: W\to Y$ is locally acyclic relative to $h^* K$. 

\end{definition}
The existence and uniqueness of $SS(K)$ is \cite[Theorem 1.3]{Beilinson}, which also proves that if $X$ has dimension $n$ then $SS(K)$ has dimension $n$ as well.


\begin{definition}
\source{bounds-stalks-perverse-sheaves-characteristic-p-shende}
\notready\cite[Definition 7.1(1)]{saito1} Let $X$ be a smooth scheme of dimension $n$ over $k$ and let $C \subseteq T^* X$ be a closed conical subset of the cotangent bundle with each irreducible component of dimension $n$. Let $W$ be a smooth scheme of dimension $m$ over $k$ and let $h: W \to X$ be a morphism over $k$.

We say that $h$ is \emph{properly $C$-transversal} if it is $C$-transversal and each irreducible component of $h^* C$ has dimension $m$. \end{definition} 

\begin{definition}
\source{bounds-stalks-perverse-sheaves-characteristic-p-shende}
\notready\cite[Definition 7.1(2)]{saito1} Let $X$ be a smooth scheme of dimension $n$ over $k$ and let $A$ be an algebraic cycle of codimension $n$ on $\subseteq T^* X$ whose support $C$ is a closed conical subset of the cotangent bundle (necessarily of dimension $n$).

 Let $W$ be a smooth scheme of dimension $m$ over $k$ and let $h: W \to X$ be a properly $C$-transversal morphism over $k$.
 
 We say that $h^{!} A$ is $(-1)^{n-m}$ times the pushforward along $dh:  W \times_X T^* X \to T^* W$ of the pullback along $h:  W \times_X T^* X \to T^* X$ of $A$, with the pullback and pushforward in the sense of intersection theory.  \end{definition}

Here the pushforward in the sense of intersection theory is well-defined because, by \cite[Lemma 3.1]{saito1}, $dh$ is finite when restricted to (the induced reduced subscheme structure) on $h^* C$, i.e finite when restricted to the support of $h^* A$. 


\begin{definition}
\source{bounds-stalks-perverse-sheaves-characteristic-p-shende}
\notready \cite[Definition 5.3(1)]{saito1} Let $X$ be a smooth scheme of dimension $n$ over $k$  and let $C\subseteq T^* X$ be a closed conical subset of the cotangent bundle. Let $Y$ be a smooth curve over $k$ and $f: X\to Y$ a morphism over $k$. 

We say a closed point $x \in X$ is at most an \emph{isolated $C$-characteristic point} of $f$ if $f$ is $C$-transversal when restricted to some open neighborhood of $x$ in $X$, minus $x$.  We say that $x\in X$ is an \emph{isolated $C$-characteristic point} of $f$ if this holds, but $f$ is not $C$-transversal when restricted to any open neighborhood of $X$. \end{definition}

\begin{definition}
\source{bounds-stalks-perverse-sheaves-characteristic-p-shende}
\notready For $V$ a representation of the Galois group of a local field over $\mathbb F_\ell$ (or a continuous $\ell$-adic representation), we define $\operatorname{dimtot} V$ to be the dimension of $V$ plus the Swan conductor of $V$. For a complex $W$ of such representations, we define $\operatorname{dimtot}W $ to be the alternating sum $\sum_i (-1)^i \operatorname{dimtot} \mathcal H^i(W)$ of the total dimensions of its cohomology objects. \end{definition}

\begin{definition}
\source{bounds-stalks-perverse-sheaves-characteristic-p-shende}
\notready\cite[Definition 5.10]{saito1} Let $X$ be a smooth scheme of dimension $n$ over $k$ and $K$ an object of $D^b_c(X, \mathbb F_\ell)$. Let the \emph{characteristic cycle of $K$}, $CC(K)$ , be the unique $\mathbb Z$-linear combination of irreducible components of $SS(K)$ such that for every \'{e}tale morphism $j: W \to X$, every morphism $f: W\to Y$ to a smooth curve and every at most isolated $h^\circ SS(\mathcal F)$-characteristic point $u \in W$ of $f$, we have

\[ - \operatorname{dimtot} \left( R \Phi_f(j^* K) \right)_u =  (j^* CC(K), (df)^*\omega )_{T^*W,u} \] where $\omega$ is a meromorphic one-form on $Y$ with no zero or pole at $f(u)$.

\end{definition}

Here the notation $(,)_{T*W ,u}$ denotes the intersection number in $T^* W$ at the point $u$.

The existence and uniqueness is \cite[Theorem 5.9]{saito1}, except for the fact that the coefficients lie in $\mathbb Z$ and not $\mathbb Z[1/p]$, which is \cite[Theorem 5.18]{saito1} and is due to Beilinson, based on a suggestion by Deligne.
